\section{Introduction}

Transformers can learn mechanisms for relational composition \citep{wattenberg2024relational}, by which multiple feature vectors are composed into a single fixed-length vector to represent complex relationships. 
For example, \citet{farrell2025order} train a transformer on an ordered bigram 
copying task that learns a ``relative magnitude-based'' composition mechanism, 
representing both $(x, y)$ and $(y,x)$ as linear combinations of the input 
embeddings with coefficients defined by the attention pattern.
This suggests that sparse autoencoders (SAEs) \citep{bricken2023monosemanticity, huben2024sparse} and their variants \citep{gao2024scaling, bussmann2024batchtopk, rajamanoharan2024jumping, bussmann2025learning, karvonen2025saebench, leask2025inference} may learn a sparse solution with latents corresponding to the input embeddings, and graded activations corresponding to ordering.
This challenges common interpretations of SAE latents as binary feature presence indicators \cite{paulo2025automatically,quirke2025binary}.
This challenge persists even if SAEs learn the correct dictionary, which separates this from previous work on SAE reliability \citep{chanin2024absorption,leask2025sparse}.

% For example, \citet{wattenberg2024relational} propose additive matrix binding. This could represent bigram $(x,y)$ as $r = Ax + By$ to preserve bigram order, using matrices to map each element onto role-specific subspaces.
% which could recover ``echo features'' $Ax, Ay, Bx, By$ rather than the actual features $x,y$.

% SAEs work by projecting dense superposed model activations into a latent space with much higher dimensionality and a sparsity penalty. In theory, this encourages the SAE to learn latents corresponding to underlying model feature directions, which recovers a dictionary of independently interpretable model features.

% SAEs have been successfully used in the past \citep{marks2025auditing, movva2025sparse}, but there are uncertainties over whether they learn robust dictionaries. For example, \citet{chanin2024absorption} find that using sparsity as a proxy can cause feature absorption and \citet{leask2025sparse} find that SAEs do not find a unique and complete dictionary of atomic features.

However, \citet{farrell2025order} did not study the impact on SAEs.
In this paper, we train a suite of BatchTopK \cite{bussmann2024batchtopk}, JumpReLU \cite{rajamanoharan2024jumping} and Matryoshka \cite{bussmann2025learning} SAEs on the toy model from \citet{farrell2025order} and find that they reliably learn latents with graded activations corresponding to different bigram orderings.
Our paper is structured as follows. Section \ref{s:method} introduces the base transformer and the task it is trained on, our SAE training setup and the steering protocol used to test for graded activations.
In Section \ref{s:results}, we analyse a specific BatchTopK SAE and find two latents with graded activations that can be steered to swap the model output from $(x,y)$ to $(y,x)$. They can be identified using their activation magnitudes, which correlate with the difference between transformer attention weights, $\Delta \alpha$.
Across a broad SAE sweep, high-performing SAEs tend to contain latents whose activations correlate with $\Delta \alpha$.
For a selected set of six high-performing SAEs across BatchTopK, JumpReLU and Matryoshka architectures, steering these latents swaps between 25-73\% of the transformer's original outputs.

Our paper is intended to provide a foundation for further research into graded latent activations in real-world pretrained SAEs such as the Gemma Scope releases \cite{lieberum2024gemma, mcdougall2025gemmascope2}, which use JumpReLU and Matryoshka SAEs.
By providing a case study with automatic detection, we aim to foster this work.
